% Original analytic teaching diagram, 2026-11-12.
% Order in which the computable information of a target is checked when choosing a sampler:
% quantile function -> inverse transform; decomposition into mixture/conditional steps -> sequential generation;
% envelope with finite constant M -> rejection; unnormalised density ratios -> MCMC.  Each method box lists the
% quantity to compare once the method is available.
\definecolor{VB}{HTML}{7188A3}\definecolor{VBF}{HTML}{E5EAF0}
\definecolor{VG}{HTML}{829C8B}\definecolor{VGF}{HTML}{E8EEE9}
\definecolor{VY}{HTML}{B4A16F}\definecolor{VYF}{HTML}{F1ECDD}
\definecolor{VR}{HTML}{A7787D}\definecolor{VRF}{HTML}{F0E5E6}
\definecolor{VP}{HTML}{9588A5}\definecolor{VPF}{HTML}{ECE8F0}
\definecolor{VI}{HTML}{3B4252}\definecolor{VM}{HTML}{6B7280}\definecolor{VL}{HTML}{4C4D4F}
\begin{tikzpicture}[x=1mm,y=1mm,>=Stealth,font=\figurefont\fontsize{8}{10}\selectfont,
lab/.style={text=VM,font=\figurefont\fontsize{7.5}{9}\selectfont},
q/.style={draw=VL,fill=white,line width=.6pt,rounded corners=1mm,align=center,text=VI,text width=36mm,minimum height=8mm,inner sep=1mm},
m/.style={draw=VB,fill=VBF,line width=.6pt,rounded corners=1mm,align=center,text=VI,text width=20mm,minimum height=8mm,inner sep=1mm},
arr/.style={->,draw=VL,line width=.6pt}]
\node[q] (q1) at (0,48) {能计算分位数函数？};
\node[q] (q2) at (0,34) {能拆成混合成分或\\[-.5mm]逐个坐标的条件分布？};
\node[q] (q3) at (0,20) {有可采样的包络，\\[-.5mm]常数 $M$ 有限且不大？};
\node[q] (q4) at (0,6) {能逐点计算\\[-.5mm]未归一化密度？};
\node[m] (m1) at (40,48) {逆变换};
\node[m] (m2) at (40,34) {分步生成};
\node[m] (m3) at (40,20) {拒绝采样};
\node[m,draw=VG,fill=VGF] (m4) at (40,6) {MCMC};
\node[lab,anchor=west,align=left] at (52,48) {数值反演的精度与成本};
\node[lab,anchor=west,align=left] at (52,34) {每步是否都可精确抽取};
\node[lab,anchor=west,align=left] at (52,20) {每个样本约 $M$ 个候选};
\node[lab,anchor=west,align=left] at (52,6) {初值、探索范围、$\tau_{f}$};
\foreach \i in {1,2,3,4}{\draw[arr] (q\i) -- node[above,lab]{是} (m\i);}
\foreach \i/\j in {1/2,2/3,3/4}{\draw[arr] (q\i) -- node[right,lab]{否} (q\j);}
\draw[arr,draw=VR] (q4.south) -- node[right,lab]{否} ++(0,-6) node[below,text=VR]{重新表述计算任务};
\end{tikzpicture}
